\documentclass[11pt,a4paper]{article}
\usepackage{fontspec,amsmath,amssymb,amsthm,microtype}
\setmainfont{lmroman10-regular.otf}[BoldFont=lmroman10-bold.otf,ItalicFont=lmroman10-italic.otf,BoldItalicFont=lmroman10-bolditalic.otf]
\setmonofont{lmmono10-regular.otf}[Scale=MatchLowercase]
\usepackage[margin=25mm]{geometry}
\usepackage[unicode,hidelinks]{hyperref}
\setlength{\emergencystretch}{3em}
\setlength{\parskip}{0.4em}
\setlength{\parindent}{0pt}
\allowdisplaybreaks[2]
\raggedbottom
\title{\texorpdfstring{Gaussian displacement and frozen-\allowbreak{}gradient error energy}{Gaussian displacement and frozen-gradient error energy}}
\author{LangevinLean mathematical exposition}
\date{}
\begin{document}\maketitle
\noindent Source module: \texttt{Unadjusted\allowbreak{}Langevin/\allowbreak{}LSI/\allowbreak{}Gaussian\allowbreak{}Drift.\allowbreak{}lean}. The accompanying HTML pairs this mathematical exposition with the exact Lean statements and records the scope of the formalization.


For \(Y_t=X-t\nabla U(X)+\sqrt{2t}Z\), independent Gaussian noise makes the cross term vanish, giving \(\mathbb E\|Y_t-X\|^2=t^2\mathbb E\|\nabla U(X)\|^2+2td\). Gradient Lipschitzness bounds the corresponding frozen-gradient error by \(L^2\) times this moment.


\section{\texorpdfstring{integral norm sq std\allowbreak{}Gaussian}{integral norm sq stdGaussian}}


Choose a finite orthonormal basis and use \(\|Z\|^2=\sum_i\langle e_i,Z\rangle^2\). Each standard Gaussian coordinate has second moment one, proved using its dual variance formula and zero mean. Gaussian \(L^2\) membership justifies integrating the finite sum, giving \(\mathbb E\|Z\|^2=d\), also when \(d=0\).



{\raggedright\noindent Lean declarations: \texttt{integral\_\allowbreak{}norm\_\allowbreak{}sq\_\allowbreak{}std\allowbreak{}Gaussian}.\par}

\section{\texorpdfstring{integral norm sq affine std\allowbreak{}Gaussian}{integral norm sq affine stdGaussian}}


Expand \(\|ag(X)+bZ\|^2=a^2\|g(X)\|^2+2ab\langle g(X),Z\rangle+b^2\|Z\|^2\). The \(L^2\) hypothesis on \(g\) and Gaussian \(L^2\) membership prove all three terms integrable on the independent product.



The product bilinear-integral theorem factors the cross expectation as \(\langle\mathbb Eg(X),\mathbb EZ\rangle=0\). The product marginals identify the first energy, and \texttt{integral\_\allowbreak{}norm\_\allowbreak{}sq\_\allowbreak{}std\allowbreak{}Gaussian} identifies the noise energy. Hence the integral is \(a^2\mathbb E\|g\|^2+b^2d\).



{\raggedright\noindent Lean declarations: \texttt{integral\_\allowbreak{}norm\_\allowbreak{}sq\_\allowbreak{}affine\_\allowbreak{}std\allowbreak{}Gaussian}.\par}

\section{\texorpdfstring{mala\allowbreak{}Proposal sub initial}{malaProposal sub initial}}


Subtract \(x\) from \(T_t(x,z)=x-t\nabla U(x)+\sqrt{2t}z\). Additive cancellation leaves \(-t\nabla U(x)+\sqrt{2t}z\). This algebraic identity shows why a state second moment is unnecessary when only displacement energy is estimated.



{\raggedright\noindent Lean declarations: \texttt{mala\allowbreak{}Proposal\_\allowbreak{}sub\_\allowbreak{}initial}.\par}

\section{\texorpdfstring{frozen displacement second moment}{frozen displacement second moment}}


Apply the affine Gaussian energy formula with \(a=-t\), \(b=\sqrt{2t}\), and \(g=\nabla U\). For \(t\ge0\), \(b^2=2t\), yielding \(\mathbb E\|T_t-X\|^2=t^2\mathbb E\|\nabla U(X)\|^2+2td\). Independence and centering have already discharged the cross term.



{\raggedright\noindent Lean declarations: \texttt{frozen\_\allowbreak{}displacement\_\allowbreak{}second\_\allowbreak{}moment}.\par}

\section{\texorpdfstring{frozen gradient error energy}{frozen gradient error energy}}


Lipschitzness gives \(\|\nabla U(T_t)-\nabla U(X)\|^2\le L^2\|T_t-X\|^2\) pointwise. Convert integrability of the squared gradient to \(L^2\) membership, so the displacement has an integrable squared norm.



Continuity of the proposal and gradient proves measurability of the squared discrepancy. Domination by the integrable displacement energy proves its integrability; \texttt{integral\_\allowbreak{}mono} and the exact displacement moment then give \(L^2(t^2G_0+2td)\). This is a joint-law estimate; a state-space conditional drift is constructed separately.



{\raggedright\noindent Lean declarations: \texttt{frozen\_\allowbreak{}gradient\_\allowbreak{}error\_\allowbreak{}energy}.\par}

\section{\texorpdfstring{frozen displacement mem\allowbreak{}Lp}{frozen displacement memLp}}


Pull the square-integrable initial gradient to the first product coordinate and the square-integrable Gaussian identity to the second. Their scalar multiples and sum lie in \(L^2\). The displacement identity rewrites this sum as \(T_t(X,Z)-X\).



{\raggedright\noindent Lean declarations: \texttt{frozen\_\allowbreak{}displacement\_\allowbreak{}mem\allowbreak{}Lp}.\par}

\end{document}
